% Generated draft: requires mathematical and source review.
\documentclass[11pt]{article}
\usepackage[margin=25mm]{geometry}
\usepackage{fontspec}
\setmainfont{Latin Modern Roman}
\usepackage{amsmath,amssymb,mathtools,unicode-math}
\setmathfont{Latin Modern Math}
\usepackage{enumitem,needspace,xurl,hyperref}
\setlength{\parindent}{0pt}
\setlength{\parskip}{5pt}
\setlength{\emergencystretch}{3em}
\newcommand{\N}{\mathbb{N}}
\newcommand{\Z}{\mathbb{Z}}
\newcommand{\Q}{\mathbb{Q}}
\newcommand{\R}{\mathbb{R}}
\newcommand{\C}{\mathbb{C}}
\newcommand{\F}{\mathbb{F}}
\newcommand{\A}{\mathbb{A}}
\newcommand{\PP}{\mathbb{P}}
\newcommand{\E}{\mathbb{E}}
\newcommand{\eps}{\varepsilon}
\newcommand{\dd}{\,\mathrm{d}}
\begin{document}
\section*{UnsolvedMath ID: 800813}
\textbf{Source problem code:} catalogue entry
\subsection{Definitions and mathematical statement}
Write \(\mathcal G\) for the class called \emph{GM-flows} in the supplied source record. The record does not define this class or specify its governing equation, spatial domain, geometric structure, boundary conditions, regularity, or time interval. These are unresolved parts of the formulation, rather than hypotheses that may be freely chosen.

For a flow \(u\in\mathcal G\), write \(M_u\) for its spatial domain and \(I_u\) for its time interval. Whenever the source's vorticity expression is a well-defined scalar, set
\[
\omega_u(x,t)=\operatorname{curl}u(x,t),\qquad
L_{u,t,c}=\{x\in M_u:\omega_u(x,t)=c\}.
\]
Here \(t\in I_u\), and \(c\) ranges over scalar vorticity values. The scalar interpretation is suggested by the source's reference to level lines, but the dimension and the definition of curl are not supplied.

\textbf{Source assertion to be proved.} For every GM-flow \(u(x,t)\), the level lines
\[
\operatorname{curl}u(x,t)=\mathrm{const}
\]
are real-analytic. The intended universal quantification is over every admissible flow, time, and vorticity level; the admissibility conditions and any exceptions are absent from the supplied record.

For reference, in a specified real-analytic ambient manifold, a regular real-analytic curve is a one-dimensional embedded submanifold locally parametrized by a real-analytic map \(\gamma:(-\delta,\delta)\to M_u\), with \(\delta>0\) and \(\gamma'\ne0\). The source does not specify whether its assertion concerns only regular level curves, entire level sets including critical points, or individual branches of singular level sets. It also does not specify the ambient analytic structure or treatment of boundary points. Consequently, this definition records one conventional meaning of real-analytic curve without selecting it as the intended interpretation. A fully quantified theorem requires the missing GM-flow definition and these level-set qualifications.
\subsection{Short English statement}
The supplied problem asks whether every vorticity level line of every GM-flow is real-analytic. The definition of GM-flow and the treatment of critical levels and boundaries are missing, so the exact theorem remains unresolved. The registry labels the problem open; this is supplied status information, not an independently verified assessment.
\subsection{Sources}
\begin{itemize}
\item Local UnsolvedMath registry, permanent ID \(800813\), exact source problem code \texttt{catalogue entry}, title \emph{Geometric Hydrodynamics in Open Problems --- Problem 12}. The record attributes the problem to Khesin, Misiolek, Modin, with proposed year \(2022\), and labels its status \texttt{open}.
\item Supplied reference: Khesin, Misiolek, Modin, \emph{Geometric Hydrodynamics in Open Problems} (2022), Problem 12, \url{https://arxiv.org/abs/2205.01143}. The background records access on 2026-10-16 and extraction by \texttt{arxiv\_latex}; these are supplied provenance claims, not a claim of independent access or verification in this draft. The supplied excerpt contains the problem statement but no definition of GM-flow.
\end{itemize}
\end{document}
